\chapter{Complex Numbers and Euler's Formula}
\label{ch:ch16-complex-euler}

\section{Extending the Number System}

Chapter~15 studied amplitwists, $A = sR(\theta)$, and found that they
compose according to the rule $(s_1,\theta_1)\circ(s_2,\theta_2) =
(s_1s_2,\theta_1+\theta_2)$, a rule that combines two different
arithmetic operations, multiplication and addition, in a single formula.
The purpose of this chapter is to show that an arithmetic answering to
exactly this description already exists: complex numbers provide an
algebraic language for amplitwists, and multiplying complex numbers turns
out to perform exactly the same operation as composing amplitwists.
Before that connection can be established, the number system itself must
first be extended.

The natural numbers arise from counting. The integers extend the natural
numbers by introducing negative quantities, allowing subtraction to be
performed universally. The rational numbers extend the integers by
introducing fractions, allowing division by any nonzero number. The real
numbers extend the rationals to accommodate measurement and continuity.
Each extension solves equations that could not previously be solved: $x+3
= 1$ has no solution among the natural numbers but has the integer
solution $x = -2$; $2x = 1$ has no integer solution but has the rational
solution $x = 1/2$; and the real numbers solve equations such as $x^2=2$,
whose solutions are not themselves rational. A natural question therefore
arises: does the equation $x^2+1=0$, which may be rewritten as $x^2=-1$,
have any solution at all? No real number has a square equal to a negative
quantity, so no real solution exists, and to solve this equation a new
kind of number must be introduced.

\section{The Imaginary Unit}

\begin{definition}
The \emph{imaginary unit} is the number $i$ satisfying $i^2 = -1$.
\end{definition}

\begin{historicalnote}{A Dismissive Name That Stuck}
The word \emph{imaginary} was coined by Ren\'e Descartes, the same figure
Chapter~2 credited with Cartesian coordinates, and he did not intend it
as a compliment \cite{boyer1991}. In his 1637 \emph{La G\'eom\'etrie},
Descartes used \emph{nombres imaginaires} to describe roots of equations
that had no correspondence to any measurable length or quantity,
essentially dismissing them as fictions useful for bookkeeping in an
algebraic calculation but not real numbers in any meaningful sense. The
name persisted long after mathematicians including Euler, Gauss, and
others showed that these numbers had a perfectly concrete geometric
meaning, as the points of a plane rather than points on a line, exactly
the picture this chapter develops through the complex plane and Euler's
formula. By the time the geometry was understood, the dismissive name
had already become standard usage, and no replacement has displaced it
since; the complex numbers of this chapter are, etymologically, still
named after Descartes' doubt that they were real numbers at all.
\end{historicalnote}

This definition immediately produces a repeating pattern:
\[
i^1 = i, \qquad i^2 = -1, \qquad i^3 = i^2\cdot i = -i, \qquad i^4 =
i^2\cdot i^2 = (-1)(-1) = 1.
\]
Continuing, $i^5 = i$ and $i^6 = -1$, and the cycle repeats indefinitely.

\begin{theorem}
The powers of $i$ repeat with period four.
\end{theorem}

\begin{proof}
Since $i^4 = 1$, for any positive integer $n$ write $n = 4q+r$ with $0
\le r < 4$; then $i^n = (i^4)^q i^r = 1^q i^r = i^r$, so $i^n$ depends
only on the remainder of $n$ upon division by $4$.
\end{proof}

\begin{example}
Compute $i^{23}$. Since $23 = 4(5)+3$, $i^{23} = (i^4)^5 i^3 = 1^5(-i) =
-i$.
\end{example}

\section{Complex Numbers}

\begin{definition}
A \emph{complex number} is an expression of the form $z = a+bi$, where
$a$ and $b$ are real numbers and $i$ is the imaginary unit.
\end{definition}

The number $a$ is called the \emph{real part} of $z$, written
$\operatorname{Re}(z) = a$, and $b$ is called the \emph{imaginary part},
written $\operatorname{Im}(z) = b$. Addition is performed componentwise,
$(a+bi)+(c+di) = (a+c)+(b+d)i$, while multiplication follows the
distributive law together with $i^2=-1$.

\begin{example}
Compute $(2+3i)(1+4i)$. Expanding gives $2+8i+3i+12i^2$; since $i^2=-1$,
this becomes $2+11i-12 = -10+11i$.
\end{example}

The algebra looks unfamiliar at first, but the next section reveals its
geometric meaning.

\section{The Complex Plane}

Complex numbers become much easier to understand once they are
represented geometrically. Recall from Chapter~2 that every ordered pair
$(a,b)$ corresponds to a point in the Cartesian plane; since every
complex number has the form $z=a+bi$, it is natural to associate $z$ with
the point $(a,b)$. This interpretation produces the \emph{complex plane},
in which the horizontal axis represents real parts and the vertical axis
represents imaginary parts:
\[
a+bi \longleftrightarrow (a,b).
\]
Thus $3+2i$ corresponds to the point $(3,2)$, while $-1+4i$ corresponds
to $(-1,4)$. Because complex numbers correspond to points, their distance
from the origin can be measured exactly as it was for vectors in
Chapter~2.

\begin{definition}
The \emph{modulus} of a complex number $z=a+bi$ is $|z| = \sqrt{a^2+b^2}$,
the length of the segment from the origin to the point representing $z$.
\end{definition}

\begin{example}
Find the modulus of $z=3+4i$. Using the definition, $|z| =
\sqrt{3^2+4^2} = \sqrt{25} = 5$, so $z$ lies five units from the origin.
\end{example}

This is exactly the distance formula introduced earlier in the text; no
new geometry has been created, only familiar geometry reinterpreted in a
new algebraic setting.

\section{Multiplication as Transformation}

Addition of complex numbers behaves much like vector addition.
Multiplication behaves differently, and to investigate its geometric
meaning, consider multiplying by $i$. For $z=a+bi$,
\[
iz = i(a+bi) = ai + bi^2 = -b+ai,
\]
so in coordinate form $(a,b) \mapsto (-b,a)$.

\begin{example}
The complex number $1$ corresponds to $(1,0)$. Multiplying by $i$ gives
$i$, corresponding to $(0,1)$: the point has rotated through $90^\circ$.
\end{example}

\begin{example}
The complex number $2+i$ corresponds to $(2,1)$. Multiplying by $i$ gives
$-1+2i$, corresponding to $(-1,2)$: again the point has rotated through
$90^\circ$.
\end{example}

\begin{theorem}
Multiplication by $i$ rotates every point in the complex plane through
$90^\circ$ counterclockwise.
\end{theorem}

\begin{proof}
Multiplication by $i$ sends $(a,b) \mapsto (-b,a)$. From Chapter~8, the
rotation matrix
\[
R(90^\circ) = \begin{pmatrix} 0 & -1\\ 1 & 0 \end{pmatrix}
\]
produces exactly the same transformation, since $\begin{pmatrix} 0 &
-1\\ 1 & 0 \end{pmatrix}\begin{pmatrix} a\\ b \end{pmatrix} =
\begin{pmatrix} -b\\ a \end{pmatrix}$. Therefore multiplication by $i$ is
a rotation through $90^\circ$.
\end{proof}

This discovery is remarkable: multiplication is no longer merely
arithmetic, but has become geometry. Multiplication by $-1$ sends
$(a,b)\mapsto(-a,-b)$, a rotation through $180^\circ$; multiplication by
$2i$ rotates through $90^\circ$ and doubles all distances from the
origin; and multiplication by $1+i$ appears to rotate and scale
simultaneously. A pattern is emerging: multiplication seems to combine
scaling and rotation into a single operation. But scaling and rotation
combined were precisely the amplitwists of Chapter~15, and the next
section makes this connection exact.

\section{Polar Form}

Chapter~14 introduced polar coordinates: every point in the plane may be
described by a distance $r$ from the origin and an angle $\theta$
measured from the positive $x$-axis, related by $x = r\cos\theta$ and $y
= r\sin\theta$. Substituting these into $z = x+iy$ gives
\[
z = r\cos\theta + ir\sin\theta,
\]
and factoring out $r$,
\[
z = r(\cos\theta + i\sin\theta).
\]

\begin{definition}
The expression $z = r(\cos\theta+i\sin\theta)$ is called the \emph{polar
form} of $z$, where $r = |z|$ is the modulus and $\theta$ is the angle
made with the positive real axis.
\end{definition}

\begin{example}
Convert $1+i$ to polar form. The modulus is $r = \sqrt{1^2+1^2} =
\sqrt2$, and since the point lies on the line $y=x$ in the first
quadrant, $\theta = 45^\circ$. Therefore $1+i = \sqrt2(\cos45^\circ +
i\sin45^\circ)$.
\end{example}

Whereas rectangular form is ideal for addition and subtraction, polar
form reveals the geometric meaning of multiplication, as the next section
shows.

\section{Multiplication in Polar Form}

The connection between complex numbers and amplitwists is revealed by
multiplying two complex numbers written in polar form. Let $z_1 =
r_1(\cos\theta_1+i\sin\theta_1)$ and $z_2 =
r_2(\cos\theta_2+i\sin\theta_2)$. Multiplying gives
\[
z_1z_2 = r_1r_2(\cos\theta_1+i\sin\theta_1)(\cos\theta_2+i\sin\theta_2),
\]
and expanding the product,
\[
z_1z_2 = r_1r_2\Big[(\cos\theta_1\cos\theta_2 - \sin\theta_1\sin\theta_2)
+ i(\sin\theta_1\cos\theta_2 + \cos\theta_1\sin\theta_2)\Big].
\]
The real part contains $\cos\theta_1\cos\theta_2 -
\sin\theta_1\sin\theta_2$, which Chapter~8 identified as
$\cos(\theta_1+\theta_2)$, while the imaginary part contains
$\sin\theta_1\cos\theta_2 + \cos\theta_1\sin\theta_2 =
\sin(\theta_1+\theta_2)$. Substituting these identities yields the
central theorem of the chapter.

\begin{theorem}
If $z_1 = r_1(\cos\theta_1+i\sin\theta_1)$ and $z_2 =
r_2(\cos\theta_2+i\sin\theta_2)$, then
\[
z_1z_2 = r_1r_2\Big(\cos(\theta_1+\theta_2) + i\sin(\theta_1+\theta_2)\Big).
\]
\end{theorem}

\begin{corollary}
When complex numbers are multiplied, moduli multiply and angles add.
\end{corollary}

This is exactly the composition law of Chapter~15. Beside the amplitwist
rule $(s_1,\theta_1)\circ(s_2,\theta_2) = (s_1s_2,\theta_1+\theta_2)$,
place the multiplication theorem just proved, $(r_1,\theta_1)\cdot(r_2,
\theta_2) = (r_1r_2,\theta_1+\theta_2)$: the two rules are identical.
Complex multiplication is amplitwist composition. Every nonzero complex
number determines an amplitwist, and multiplying complex numbers
composes the corresponding transformations; the geometry of Chapter~15
and the algebra of this chapter describe the same mathematical structure.

\begin{example}
Multiply $2(\cos30^\circ+i\sin30^\circ)$ and
$3(\cos40^\circ+i\sin40^\circ)$. The theorem gives $z_1z_2 =
6(\cos70^\circ+i\sin70^\circ)$: the modulus has become $2\cdot3=6$, while
the angle has become $30^\circ+40^\circ=70^\circ$.
\end{example}

\section{Euler's Formula}

The expression $\cos\theta+i\sin\theta$ occurs so frequently that a
special notation is introduced for it.

\begin{definition}[Euler's Formula]
For every real number $\theta$, $e^{i\theta} = \cos\theta+i\sin\theta$.
\end{definition}

Using this notation, a complex number in polar form becomes $z =
re^{i\theta}$, and the multiplication theorem of the previous section
takes its simplest possible form:
\[
e^{i\theta_1}e^{i\theta_2} = e^{i(\theta_1+\theta_2)}, \qquad\text{so
that}\qquad (re^{i\theta})(se^{i\phi}) = rse^{i(\theta+\phi)}.
\]
Angle addition has become ordinary exponent addition. This chapter treats
Euler's formula as a definition motivated by the compactness it lends to
the multiplication law already proved; a derivation of the same formula
from the power series for $e^x$, $\sin x$, and $\cos x$ belongs to a
calculus course and is not needed here to use the formula correctly.
Instead of carrying trigonometric expressions through every calculation,
this notation allows complex arithmetic to be handled directly through
exponentials, with multiplication automatically scaling by $rs$ and
rotating by $\theta+\phi$ in a single step.

\section{Euler's Identity}

A celebrated special case occurs when $\theta = \pi$. Substituting into
Euler's formula gives $e^{i\pi} = \cos\pi + i\sin\pi$, and since $\cos\pi
= -1$ and $\sin\pi = 0$, this becomes $e^{i\pi} = -1$, or equivalently:

\begin{theorem}[Euler's Identity]
$e^{i\pi} + 1 = 0$.
\end{theorem}

This equation connects five of the most important constants in
mathematics --- $0$, $1$, $\pi$, $e$, and $i$ --- in a single relation,
and is often regarded as one of the most elegant formulas in mathematics
for exactly that reason.

\section{Repeated Multiplication and De Moivre's Theorem}

Chapter~15 showed that repeated amplitwist composition satisfies $A^n =
s^nR(n\theta)$. Complex numbers obey an identical law. Let $z =
re^{i\theta}$; then $z^2 = (re^{i\theta})(re^{i\theta}) = r^2e^{i2\theta}$,
and similarly $z^3 = r^3e^{i3\theta}$, suggesting the general formula $z^n
= r^ne^{in\theta}$.

\begin{theorem}[De Moivre's Theorem]
For every positive integer $n$, $(re^{i\theta})^n = r^ne^{in\theta}$.
Equivalently, $(\cos\theta+i\sin\theta)^n = \cos(n\theta)+i\sin(n\theta)$.
\end{theorem}

\begin{proof}
Using the multiplication law for complex numbers in polar form
repeatedly, each additional factor of $re^{i\theta}$ contributes one more
factor of $r$ to the modulus and adds one more $\theta$ to the angle:
$(re^{i\theta})^2 = r^2e^{i2\theta}$, $(re^{i\theta})^3 = r^3e^{i3\theta}$,
and continuing this process $n$ times yields $(re^{i\theta})^n =
r^ne^{in\theta}$. Setting $r=1$ and applying Euler's formula gives
$(e^{i\theta})^n = e^{in\theta}$, which is
$(\cos\theta+i\sin\theta)^n = \cos(n\theta)+i\sin(n\theta)$.
\end{proof}

The similarity to Chapter~15 is striking: $A^n = s^nR(n\theta)$ and
$(re^{i\theta})^n = r^ne^{in\theta}$ are the same statement written in
two different languages, the first describing repeated amplitwist
composition geometrically and the second describing repeated complex
multiplication algebraically.

\begin{example}
Compute $(\cos30^\circ+i\sin30^\circ)^6$. By De Moivre's theorem, this
equals $\cos(180^\circ)+i\sin(180^\circ)$; since $\cos180^\circ=-1$ and
$\sin180^\circ=0$, the value is $-1$.
\end{example}

\begin{algorithmproc}{How to Compute Powers Using Euler's Formula}
Write the complex number in polar form $z = re^{i\theta}$, apply De
Moivre's theorem to obtain $z^n = r^ne^{in\theta}$ by raising the modulus
$r$ to the $n$-th power and multiplying the angle $\theta$ by $n$, and
convert the result back to rectangular form if required.
\end{algorithmproc}

\begin{practicenow}
Compute $\left(2(\cos45^\circ+i\sin45^\circ)\right)^4$.
\end{practicenow}

\section{Roots of Unity}

The equation $z^n=1$ asks for all complex numbers whose $n$-th power
equals $1$. Using polar form, write $1 = e^{i2\pi k}$ for any integer
$k$, and let $z = re^{i\theta}$, so that $z^n = r^ne^{in\theta}$. Setting
this equal to $e^{i2\pi k}$ requires equal moduli and equal angles, so
$r^n=1$, giving $r=1$, and $n\theta = 2\pi k$, giving $\theta = 2\pi
k/n$.

\begin{theorem}
The solutions of $z^n=1$ are $z_k = e^{2\pi ik/n}$ for $k = 0, 1, 2,
\dots, n-1$, called the \emph{$n$-th roots of unity}.
\end{theorem}

Each root lies on the unit circle, and successive roots differ by a
rotation through $2\pi/n$, so the roots form the vertices of a regular
$n$-gon.

\begin{example}
Find the cube roots of unity. Here $n=3$, so $z_k = e^{2\pi ik/3}$. For
$k=0$, $z_0=1$. For $k=1$, $z_1 = \cos120^\circ+i\sin120^\circ =
-\tfrac12+\tfrac{\sqrt3}{2}i$. For $k=2$, $z_2 = \cos240^\circ +
i\sin240^\circ = -\tfrac12-\tfrac{\sqrt3}{2}i$. These three points form
an equilateral triangle on the unit circle.
\end{example}

From the amplitwist viewpoint, the roots of unity are equally
remarkable: each root represents a rotation which, applied repeatedly $n$
times, returns exactly to the starting position.

\section{Applications}

Complex numbers appear throughout mathematics, science, and engineering.
In electrical engineering, alternating-current circuits are often
described using complex numbers whose moduli represent amplitudes and
whose angles represent phase shifts. In signal processing, oscillations
are represented by rotating complex numbers called \emph{phasors}, and
complex multiplication allows amplitude and phase changes to be handled
simultaneously. In computer graphics, rotations and scalings can be
performed through complex multiplication, making many geometric
calculations more efficient. In physics, wave phenomena are frequently
modeled using expressions involving Euler's formula, since oscillatory
motion can be represented compactly by complex exponentials. Although
these applications may appear different on the surface, they all exploit
the same fact established in this chapter: complex multiplication
combines scaling and rotation into a single operation.

\section{Looking Ahead}

This chapter has shown that the amplitwists of Chapter~15 and the
complex numbers of this chapter describe the same mathematical
structure: a complex number written as $re^{i\theta}$ represents an
amplitwist with amplitude $r$ and angle $\theta$, and multiplication
obeys the familiar law $(re^{i\theta})(se^{i\phi}) =
rse^{i(\theta+\phi)}$. Geometry has become algebra. The next chapter
begins a new phase of the book, moving from geometric transformations to
analytical questions involving limits and rates of change. Its first goal
is to establish the fundamental limit $\lim_{\theta\to0} \sin\theta/\theta
= 1$, which leads directly to the derivatives of the sine and cosine
functions and forms the bridge from trigonometry to calculus.

\section{Exercises}

\subsection*{Basic Computations}
\begin{exercise}
Compute $i^7$.
\end{exercise}
\begin{exercise}
Compute $i^{50}$.
\end{exercise}
\begin{exercise}
Simplify $(3+2i)+(4-5i)$.
\end{exercise}
\begin{exercise}
Simplify $(2+3i)(1-2i)$.
\end{exercise}
\begin{exercise}
Find the modulus of $3+4i$.
\end{exercise}

\subsection*{The Complex Plane}
\begin{exercise}
Plot the complex number $2+3i$.
\end{exercise}
\begin{exercise}
Plot the complex number $-4+i$.
\end{exercise}
\begin{exercise}
Plot the complex number $-2-2i$, and describe its position relative to
the axes.
\end{exercise}
\begin{exercise}
Explain how the complex plane extends the Cartesian plane introduced in
Chapter~2.
\end{exercise}

\subsection*{Multiplication and Geometry}
\begin{exercise}
Verify that multiplication by $i$ sends $(a,b)$ to $(-b,a)$.
\end{exercise}
\begin{exercise}
Determine the image of $3+i$ under multiplication by $i$.
\end{exercise}
\begin{exercise}
Determine the image of $2-4i$ under multiplication by $i$.
\end{exercise}
\begin{exercise}
Describe geometrically the effect of multiplication by $-1$.
\end{exercise}
\begin{exercise}
Explain why multiplication by $i$ represents a rotation rather than a
translation.
\end{exercise}

\subsection*{Polar Form}
\begin{exercise}
Convert $1+i$ to polar form.
\end{exercise}
\begin{exercise}
Convert $-1+i$ to polar form.
\end{exercise}
\begin{exercise}
Convert $2(\cos60^\circ+i\sin60^\circ)$ to rectangular form.
\end{exercise}
\begin{exercise}
Convert $3(\cos150^\circ+i\sin150^\circ)$ to rectangular form.
\end{exercise}
\begin{exercise}
Explain why polar form is especially useful for multiplication.
\end{exercise}

\subsection*{Euler's Formula and De Moivre's Theorem}
\begin{exercise}
Use Euler's formula to evaluate $e^{i0}$.
\end{exercise}
\begin{exercise}
Use Euler's formula to evaluate $e^{i\pi/2}$.
\end{exercise}
\begin{exercise}
Use Euler's formula to evaluate $e^{i\pi}$.
\end{exercise}
\begin{exercise}
Use Euler's formula to evaluate $e^{i3\pi/2}$.
\end{exercise}
\begin{exercise}
Use Euler's formula to evaluate $e^{i2\pi}$.
\end{exercise}
\begin{exercise}
Compute $(\cos30^\circ+i\sin30^\circ)^2$.
\end{exercise}
\begin{exercise}
Compute $(\cos45^\circ+i\sin45^\circ)^4$.
\end{exercise}
\begin{exercise}
Compute $(\cos60^\circ+i\sin60^\circ)^3$.
\end{exercise}
\begin{exercise}
Compute $\left(2(\cos30^\circ+i\sin30^\circ)\right)^3$.
\end{exercise}
\begin{exercise}
Compute $\left(3(\cos20^\circ+i\sin20^\circ)\right)^4$.
\end{exercise}
\begin{exercise}
Express $(\cos\theta+i\sin\theta)^5$ using De Moivre's theorem.
\end{exercise}
\begin{exercise}
Explain geometrically why De Moivre's theorem multiplies angles by $n$.
\end{exercise}

\subsection*{Roots of Unity}
\begin{exercise}
Find all solutions of $z^2=1$.
\end{exercise}
\begin{exercise}
Find all solutions of $z^3=1$.
\end{exercise}
\begin{exercise}
Find all solutions of $z^4=1$.
\end{exercise}
\begin{exercise}
Find all solutions of $z^6=1$.
\end{exercise}
\begin{exercise}
Sketch the fourth roots of unity on the complex plane.
\end{exercise}
\begin{exercise}
Sketch the sixth roots of unity on the complex plane.
\end{exercise}
\begin{exercise}
Show that the cube roots of unity form an equilateral triangle.
\end{exercise}
\begin{exercise}
Show that the fourth roots of unity form a square.
\end{exercise}
\begin{exercise}
Determine the angle separating adjacent eighth roots of unity.
\end{exercise}
\begin{exercise}
Explain why the $n$-th roots of unity always form a regular polygon.
\end{exercise}

\subsection*{Connections with Amplitwists}
\begin{exercise}
Write the complex number $2e^{i45^\circ}$ in rectangular form.
\end{exercise}
\begin{exercise}
Write the complex number $3e^{i120^\circ}$ in rectangular form.
\end{exercise}
\begin{exercise}
Multiply $2e^{i30^\circ}$ and $3e^{i40^\circ}$.
\end{exercise}
\begin{exercise}
Show that multiplying by $e^{i\theta}$ leaves the modulus of a complex
number unchanged.
\end{exercise}
\begin{exercise}
Show that multiplying by $re^{i0}$ scales a complex number without
changing its angle.
\end{exercise}
\begin{exercise}
Explain why $(re^{i\theta})(se^{i\phi}) = rse^{i(\theta+\phi)}$ is the
amplitwist composition law of Chapter~15 written algebraically.
\end{exercise}
\begin{exercise}
Describe geometrically the transformation represented by multiplication
by $1+i$.
\end{exercise}
\begin{exercise}
Determine the modulus and angle of $1+i$.
\end{exercise}
\begin{exercise}
Show that every nonzero complex number determines a unique amplitwist.
\end{exercise}
\begin{exercise}
Explain why Chapter~15 can be viewed as a geometric introduction to
complex multiplication.
\end{exercise}

\subsection*{Proofs and Theory}
\begin{exercise}
Prove that the powers of $i$ repeat with period four.
\end{exercise}
\begin{exercise}
Prove that the product of two nonzero complex numbers is nonzero.
\end{exercise}
\begin{exercise}
Show that multiplication by $i$ preserves distance from the origin.
\end{exercise}
\begin{exercise}
Prove that multiplication by $i$ represents a $90^\circ$ rotation.
\end{exercise}
\begin{exercise}
Prove De Moivre's theorem by mathematical induction.
\end{exercise}
\begin{exercise}
Show directly from Euler's formula that $e^{i\theta_1}e^{i\theta_2} =
e^{i(\theta_1+\theta_2)}$.
\end{exercise}
\begin{exercise}
Show that every nonzero complex number can be written in polar form.
\end{exercise}
\begin{exercise}
Show that the modulus of a product satisfies $|z_1z_2| = |z_1||z_2|$.
\end{exercise}
\begin{exercise}
Show that the angle of a product equals the sum of the angles of the
factors.
\end{exercise}
\begin{exercise}
Explain why complex multiplication naturally combines scaling and
rotation.
\end{exercise}

\subsection*{Challenge Problems}
\begin{exercise}
Determine all solutions of $z^5=1$.
\end{exercise}
\begin{exercise}
Determine all solutions of $z^8=1$.
\end{exercise}
\begin{exercise}
Show that the product of all $n$-th roots of unity is $(-1)^{n+1}$.
\end{exercise}
\begin{exercise}
Suppose $z=e^{i\theta}$ is an $n$-th root of unity with $z\neq1$. Show
that $1+z+z^2+\cdots+z^{n-1}=0$.
\end{exercise}
\begin{exercise}
Prove that the $n$-th roots of unity form a group under multiplication.
\end{exercise}
\begin{exercise}
Show that repeated multiplication by a complex number of modulus $1$
traces points on a circle.
\end{exercise}
\begin{exercise}
Let $z=re^{i\theta}$. Determine all values of $r$ and $\theta$ for which
$z^2=1$.
\end{exercise}
\begin{exercise}
Determine all values of $z$ satisfying $z^3=1$, expressing each in both
polar and rectangular form.
\end{exercise}
\begin{exercise}
Show that complex multiplication is commutative by using the amplitwist
interpretation.
\end{exercise}
\begin{exercise}
Chapter~15 showed that every orientation-preserving similarity
transformation fixing the origin can be written uniquely as $A =
sR(\theta)$. Restate this result in the language of complex numbers, and
explain what it implies about every nonzero complex number.
\end{exercise}
\begin{exercise}
Chapter~15 described amplitwists geometrically, and this chapter
described complex numbers algebraically. Explain in detail why these are
not two different theories but two descriptions of the same mathematical
structure.
\end{exercise}
